\chapter[Band 40: Gruppen]{Band 40 auf einen Blick}
\noindent\textbf{Gruppen}\par\medskip
Vorausgesetzt werden die Halbgruppen und die Monoide aus Band~38.
Operation und neutrales Element gehören stets zur Strukturbezeichnung.

\begin{BandUebersicht}
\textbf{Gruppenaxiome}

\UebFormel{\begin{aligned}&\GroupAlg G\star e\ \leftrightarrow\\&\MonoidAlg G\star e\ \land\\&\forall x\in G\,\exists y\in G\\&\quad(x\star y=e\land y\star x=e).\end{aligned}}
\UebQuellen{\FormulaRefAuto{GroupDef}[def]}
&
\textbf{Inverse und Kürzung}
Für \(x,y,z\in G\):
\UebFormel{\begin{aligned}&(x\star y)^{-1}=y^{-1}\star x^{-1},\\&x\star y=x\star z\ \vdash y=z,\\&y\star x=z\star x\ \vdash y=z.\end{aligned}}
\UebQuellen{\FormulaRefAuto{GroupInverseUnique}[thm];
\FormulaRefAuto{GroupProductInverse}[thm];
\FormulaRefAuto{GroupLeftCancellation}[thm];
\FormulaRefAuto{GroupRightCancellation}[thm]} \\
\addlinespace[.8ex]

\textbf{Gruppenhomomorphismen}
Eine Abbildung zwischen Gruppen erhält die Operation; Isomorphismen sind bijektiv.
\UebFormel{f(x\star y)=f(x)\diamond f(y).}
\UebQuellen{\FormulaRefAuto{GroupHomDef}[def];
\FormulaRefAuto{GroupIsoDef}[def]}
&
\textbf{Erhaltung}

\UebFormel{\begin{aligned}&f(e_G)=e_H,\\&f(x^{-1})=f(x)^{-1}.\end{aligned}}
\UebQuellen{\FormulaRefAuto{GroupHomIdentityPreservation}[thm];
\FormulaRefAuto{GroupHomInversePreservation}[thm];
\FormulaRefAuto{GroupHomComposition}[thm]} \\
\addlinespace[.8ex]

\textbf{Einheitengruppe}

\UebFormel{\begin{aligned}&A^\times=\{u\in A\mid\MonoidUnit{A,\star,e}{u}\},\\&\MonoidAlg A\star e\ \vdash\GroupAlg{A^\times}\star e.\end{aligned}}
\UebQuellen{\FormulaRefAuto{MonoidUnitSetDef}[def];
\FormulaRefAuto{MonoidUnitSetIsGroup}[thm]}
&
\textbf{Rekonstruktion der Einheiten}
Für einen Isomorphismus \(F\) der großen Potenzhalbgruppen zweier Monoide:
\UebFormel{\begin{aligned}&F(A^\times)=B^\times,\\&F[\PowNE{A^\times}]=\PowNE{B^\times}.\end{aligned}}
\UebQuellen{\FormulaRefAuto{LargePowerSemigroupUnitGroupReconstruction}[thm]} \\
\addlinespace[.8ex]

\textbf{Gruppenpotenzhalbgruppe}

\UebFormel{\begin{aligned}&X\star_{\mathcal P}Y=\{x\star y\mid x\in X,y\in Y\},\\&\mathrm{Sing}(G)=\{\{g\}\mid g\in G\}.\end{aligned}}
\UebQuellen{\FormulaRefAuto{PowerSemigroupProductDef}[def];
\FormulaRefAuto{PowerGroupUnitsAreSingletons}[thm]}
&
\textbf{Einseitige Gruppenstarrheit}
Ist \(G\) eine Gruppe und \(B\) eine Halbgruppe, so trägt auch \(B\) eine Gruppenstruktur.
\UebFormel{\begin{aligned}&(\PowNE G,\star_{\mathcal P})\cong(\PowNE B,\diamond_{\mathcal P})\\&\qquad\vdash (G,\star)\cong(B,\diamond).\end{aligned}}
\UebQuellen{\FormulaRefAuto{GroupPowerSemigroupRigidity}[thm];
\FormulaRefAuto{PowerGroupsSingletonTransport}[thm]} \\
\addlinespace[.8ex]

\textbf{Endliche Kürzbarkeit}

\UebFormel{\CancellativeSemigroupAlg A\star,\quad\Finite A,\quad A\ne\varnothing.}
\UebQuellen{\FormulaRefAuto{CancellativeSemigroupDef}[def]}
&
\textbf{Gewinnung einer Gruppe}

\UebFormel{\exists e\in A\,\GroupAlg A\star e.}
\UebQuellen{\FormulaRefAuto{FiniteCancellativeSemigroupIsGroup}[thm]} \\
\addlinespace[.8ex]
\end{BandUebersicht}
